\documentclass[fleqn, 12pt]{article}
\usepackage{amsmath,amsthm,amssymb,amsfonts,setspace}
\setlength{\mathindent}{0pt}
\usepackage[utf8]{inputenc}
\usepackage[brazil]{babel}

\textwidth 7.0 truein
\oddsidemargin -0.25in   %left-hand edge
\evensidemargin -0.5 truein  %right-hand edge
\topmargin -0.85in      %top of paper to top of head, pulls whole unit
\textheight 9.5in

%%%% In most cases you won't need to edit anything above this line %%%%

\begin{document}
\hfill Matheus Ferreira Souza 

\hfill Lista - 004

\hfill Folha de respostas


\begin{itemize}
\item[1)] Ex001 -\\Parâmetros: $\mu = 240\,\text{ms}\,;\,\sigma=35\,\text{ms}$\\

\textbf{a)}
\begin{align*}
    P&(X<T) = 0,92\\
    X&= \mu + z\sigma\\
    &= 240 + (1,41 * 35)\\
    &= 240 +49,35 = 289,35\,\text{ms}
\end{align*}


\textbf{b)}
\begin{align*}
    Z &= \frac{200 - 240}{35}= -40/35 = - 1,14\\
    Z&= \frac{280 - 240}{35}= 40/35 = 1,14\\\\
    P&(200 \leq X \leq 280) \rightarrow P(-1,14 \leq X \leq 1,14)\\
    P&= 0,8729 - 0,1271\\
    &=0,7488 \rightarrow P = 74,88\%
\end{align*}


\textbf{c)}
\begin{align*}
    P&(X > 300) = P\left(Z = \frac{300 - 240}{35}\right)\\
    P&= 60/35 = 1,71\\\\
    P&(X > 300) \rightarrow P(Z > 1,71)\\
    &= 0,9564\\
    &=1 - 0,9564 = 0,0436\\
    P&=n * p\\
    &= 10.000 * 0,0436 = 436
\end{align*}


\textbf{d)}
\begin{align*}
    P&(X \leq x) = 0,99 \rightarrow Z_{0,99} = 2,33\\
    X &= \mu + z\sigma\\
    &= 240 + (2,33 * 35)\\
    &= 240 + 81,55 = 321,55\,\text{ms}
\end{align*}

\item[2)] Ex002 -\\

\textbf{a)}
\begin{align*}
    \text{Equipe A:}\\ 
    CV &=\frac{4}{18} * 100\\
    &=0,2222 * 100\\
    &= 22,22\%\\
    \text{Equipe B:}\\
    CV &=\frac{9}{25} * 100\\\
    &= 0,36 * 100\\
    &= 36\%
\end{align*}
A equipe A apresenta maior consistência relativa por ter um coeficiente de variação menor.\\

\textbf{b)}
\begin{align*}
    Z_A &= \frac{26 - 18}{4} = 2\\\\
    Z_B &= \frac{37 - 25}{9} = 1,33
\end{align*}
O desenvolvedor A foi o mais excepcional porque seu resultado está mais distante da média da própria equipe.\\

\textbf{c)}
\begin{align*}
    P&(X > 25) \rightarrow P\left(Z = \frac{25 - 18}{4}\right)\\
    P &= 7/4 = 1,75\\
    &= 0,9599\\
    &= 1 - 0,9599 = 0,0401
\end{align*}
A probabilidade de um integrante da equipe A produzir mais que a média da equipe B é de 4,01\%.\\

\textbf{d)}
\begin{align*}
    &\mu \pm 2\sigma \rightarrow 25\pm(2*9)\\
    &25 + 18 = 43\\
    &25 - 18 = 7
\end{align*}
Aproximadamente 95\% dos desenvolvedores entregariam entre 7 e 43 tarefas por dia.\\

\item[3)] Ex003 -\\ Parâmetros: $\mu = 14\,;\,\sigma =3$\\

\textbf{a)}
\begin{align*}
    P&(X > x) =0,02 \rightarrow P(X > x) = 0,98\\
    X &= 14 + (2,05 * 3)\\
    &= 14 + 6,15 = 20,15\\
    &\approx 21
\end{align*}
O dia se torna instável a partir de 21 falhas.\\

\textbf{b)}
\begin{align*}
    Z &= \frac{10 - 14}{3}= -4/3 = -1,33\\
    Z & = \frac{18 - 14}{3}= 4/3 = 1,33\\\\
    P&(-1,33 \leq X \leq 1,33)\\
    P &= 0,9082 - 0,0918 = 0,8164
\end{align*}
A probabilidade de ocorrerem entre 10 e 18 falhas é de 81,64\%.

\textbf{c)}
\begin{align*}
    E = 180 * 0,02 \approx 4
\end{align*}
São esperados 4 dias instáveis em um período de 180 dias.\\

\textbf{d)}
Embora a média do sistema permaneça inalterada, o aumento do desvio padrão indica maior dispersão dos dados em torno da média. Isso implica maior variabilidade no número de falhas diárias, redução de previsibilidade e aumento da probabilidade de ocorrência de valores extremos. Consequentemente, o comportamento do sistema torna-se menos estável e menos consistente.\newpage

\item[4)] Ex004 -\\ Parâmetros: $\mu = 80 \,;\, \sigma = 12$\\

\textbf{a)}
\begin{align*}
    P&(N \geq 95) \rightarrow P(X \geq 94,5)\\\\
    Z&=\frac{94,5 - 80}{12}=14,5/12 = 1,21\\\\
    P&(Z \geq 1,21) \approx 0,8869\\
    &=1 - 0,8869 = 0,1131
\end{align*}
Existe cerca de 11,31\% chance de ocorrerem pelo menos 95 erros.\\

\textbf{b)}
\begin{align*}
    P&(69,5 \leq X \leq 85,5)\\
    Z&=\frac{69,5 - 80}{12}=-10,5/12 = -0,88\\
    Z&=\frac{85,5 - 80}{12}=5,5/12 = 0,46\\\\
    P&(Z\leq-0,88)\approx0,1894\\
    P&(Z\leq0,46)\approx0,6772\\
    P&=0,6772 - 0,1894 \approx 0,4878
\end{align*}
Existe cerca de 48,78\% de chance de uma build apresentar entre 70 e 85 erros.\\

\textbf{c)}
\begin{align*}
    P&(89,5 \leq X \leq 90,5)\\
    Z&=\frac{89,5 - 80}{12}=9,5/12\approx0,79\\
    Z&=\frac{90,5 - 80}{12}=10,5/12\approx0,88\\\\
    P&(Z\leq 0,79) = 0,7852\\
    P&(Z\leq 0,88) =0,8106\\
    P &= 0,8106 - 0,7852 = 0,0254
\end{align*}
A probabilidade de uma build possuir exatamente 90 erros é de 2,54\%.\\

\textbf{d)}
A correção de continuidade melhora a aproximação normal porque a distribuição normal é contínua, enquanto muitos problemas reais, como uma quantidade de erros, são discretos. O ajuste de 0,5 unidade torna a área da curva normal mais próxima dos valores inteiros reais, aumentando a precisão da aproximação.

\item[5)] Ex005 -\\Parâmetros: $\sigma = 2,8\,\text{horas}$

\textbf{a)}
\begin{align*}
    n&=\left (\frac{1,96 * 2,8}{0,4} \right)^2\\
    n&=\left (\frac{5,48}{0,4} \right)^2 = 13,72^2 \approx 188,24\\
    n&=189
\end{align*}

\textbf{b)}
\begin{align*}
    n&=\left (\frac{2,576 * 2,80}{0,2} \right)^2\\
    n&=\left (\frac{7,2128}{0,2} \right)^2 = 36,064^2 \approx 1300,61\\
    n&=1301
\end{align*}

\textbf{c)}
O aumento drástico do tamanho da amostra ocorre porque reduzir a margem de erro pela metade quadruplica a amostra. Outro fator importante é que aumentar o nível de confiança também aumenta o valor crítico Z, ampliando o tamanho necessário da amostra.\newpage

\item[6)] Ex006 -\\Parâmetros: $Z_{\alpha/2} = 2,326\,;\,E =0,025$\\

\textbf{a)}
\begin{align*}
    n&=\frac{2,326^2 * 0,5 * 0,5}{0,025^2}\\
    &=\frac{5,410276*0,25}{0,000625}\\
    &=1,352569/0,000625\\
    &=2164,11\\\\
    n&=2165
\end{align*}

\textbf{b)}
Parâmetros: p=0,18\,;\,q=0,82\\
\begin{align*}
    n&=\frac{2,326^2*0,18*0,82}{0,025^2}\\
    &=\frac{5,410276*0,1476}{0,000625}\\
    &=0,7985567376/0,000625\\
    &=1277,69\\\\
    n&=1278
\end{align*}

\textbf{c)}
Usar p = 0,5 é considerado o cenário mais conservador porque esse valor gera a maior variabilidade possível nos resultados da pesquisa, como consequência, o cálculo exige a maior quantidade de entrevistas. Dessa forma, mesmo sem possuir uma estimativa prévia da proporção real, a amostra calculada ainda será suficiente para garantir o nível de confiança e o erro desejado.\\

\textbf{d)}
\begin{align*}
    E&=Z_{\alpha/2}\sqrt{\frac{p(1-p)}{n}} \rightarrow E = 2,326\sqrt{\frac{0,5*(1-0,5)}{600}}\\
    &=2,326\sqrt{\frac{0,25}{600}}\\
    &=2,326\sqrt{0,0004167} = 2,326*0,02041\\
    E&\approx0,0475\\\\
    E&=4,75\%
\end{align*}
\newpage
\item[7)] Ex007 -

\textbf{a)}
A amostragem estratificada é mais adequada porque os setores possuem características diferentes e variabilidades distintas. Dividir a população em estratos homogêneos reduz a variância das estimativas e melhora a precisão da média estimada da produtividade. Além disso, garante representação adequada de todos os setores da empresa.\\

\textbf{b)}
\begin{align*}
    180&*6=1080\\
    120&*8=960\\
    60&*12=720\\
    40&*5=200
\end{align*}
Os setores backend e frontend possuem a maior influência na variabilidade global, pois combinam grande quantidade de funcionários com desvios padrão relativamente elevados. O setor de dados possui alta variabilidade individual, mas menor impacto total devido ao tamanho reduzido do estrato.\\

\textbf{c)}
Na alocação ótima, os estratos recebem quantidades de amostra proporcionais ao tamanho e à variabilidade de cada setor. Assim, setores maiores e mais heterogêneos recebem mais observações, pois contribuem mais para a variância total da estimativa.\\

\textbf{d)}
A alocação uniforme distribui o mesmo número de amostras para todos os setores, ignorando diferenças de tamanho e variabilidade, o que pode gerar desperdício de observações e menor precisão. Já a alocação proporcional ajusta a amostra ao tamanho de cada setor, representando melhor a população.

\newpage
\item[8)] Ex008 - Parâmetros: $\mu=3,2\,\text{s}\,;\,\sigma=0,7\,\text{s}\,;\,p=0,22\,;\,q=0,78$

\textbf{a)}
\begin{align*}
    P&(X>5)\\
    Z&=\frac{5-3,2}{0,7} = 1,8/0,7\approx2,57\\\\
    P&(Z>2,57) \rightarrow P(Z<2,57) = 0,9949\\
    &=1-0,9949 =0,0051\\\\
    P&(X>5)=0,51\%
\end{align*}

\textbf{b)}
\begin{align*}
    &\mu\pm2\sigma \rightarrow 3,2\pm2(0,7)\\
    &\text{Limite inferior:}\;3,2-1,4=1,8\,\text{s}\\
    &\text{Limite superior:}\;3,2+1,4=4,6\,\text{s}\\\\
    &P(1,8\,\text{s}\leq X\leq4,6\,\text{s})
\end{align*}

\textbf{c)}
Parâmetros: $Z_{0,90}\approx1,28$\\
\begin{align*}
    X&=\mu+Z\sigma\\
    &=3,2+(1,28*0,7) =3,2+0,896\\
    &=4,10\\\\
    P&_{90}\approx4,1\,\text{s}
\end{align*}

\textbf{d)}
\begin{align*}
    n&=\frac{Z^2pq}{E^2} \rightarrow \frac{2,576^2*0,22*0,78}{0,015^2}\\
    &=\frac{6,635776*0,1716}{0,000225}\\
    &=1,1391/0,000225\\
    n&\approx5062,7\\\\
    n&=5063
\end{align*}

\textbf{e)}
O erro amostral impacta mais fortemente o tamanho da amostra, pois aparece no denominador da fórmula elevado ao quadrado. Assim, pequenas reduções no erro produzem grandes aumentos na amostra necessária. Já o nível de confiança altera apenas o valor crítico Z, causando impacto relativamente menor.

\end{itemize}

\end{document}